\chapter{Revisão Bibliográfica}
\label{chap3}

Este capítulo apresenta a revisão bibliográfica que fundamenta a análise desenvolvida neste trabalho, com foco nas abordagens de modelagem aplicadas a sistemas de destilação por membranas. Inicialmente, são retomados aspectos gerais do processo de destilação por membranas, de modo a situar o leitor quanto aos fenômenos físicos envolvidos e às variáveis de desempenho de interesse. Em seguida, são discutidas as principais estratégias de modelagem adotadas na literatura, incluindo modelos baseados em primeiros princípios, modelos orientados por dados e abordagens híbridas, destacando suas vantagens, limitações e contribuições para o estudo de sistemas de MD. A partir dessa análise, são identificadas lacunas que justificam a proposta metodológica adotada neste estudo.

% CITAÇÃO NECESSÁRIA: Incluir uma referência geral para a definição e os princípios da destilação por membranas, caso as referências ExperimentalMD e S2MD não sejam suficientes para sustentar este parágrafo.
A destilação por membranas é um processo de separação térmico baseado no transporte de vapor através de uma membrana hidrofóbica e microporosa. A força motriz do processo é a diferença de pressão parcial de vapor entre os dois lados da membrana, geralmente estabelecida por uma diferença de temperatura entre as correntes. Nesse processo, a fase líquida é impedida de atravessar os poros da membrana, enquanto o vapor de água migra para o lado de menor pressão parcial, onde é posteriormente condensado.

\section{Modelagem de Sistemas de Destilação por Membranas}
\label{sec:modelagem_md}

A modelagem matemática de sistemas de destilação por membranas é estudada porque permite descrever, prever e analisar o desempenho do processo sem depender exclusivamente de campanhas experimentais extensas. Em sistemas de MD, essa previsão é particularmente relevante porque o desempenho do processo envolve tanto a produtividade, associada ao fluxo de permeado, quanto parâmetros de eficiência energética, que dependem das temperaturas de saída das correntes quente e fria. Essas variáveis são determinadas por fenômenos simultâneos de transferência de calor e massa, evaporação na interface quente, transporte de vapor através da membrana e condensação no lado frio (\cite{ExperimentalMD, S2MD}). Assim, modelos matemáticos são ferramentas importantes tanto para a compreensão dos mecanismos dominantes quanto para o projeto, otimização e ampliação de escala desses sistemas.

Entretanto, a modelagem desses sistemas permanece desafiadora devido ao caráter acoplado e não linear dos fenômenos envolvidos. A resposta do processo depende de variáveis operacionais como temperatura de alimentação, temperatura de resfriamento, vazões, pressão de vácuo e salinidade, além de propriedades da membrana e de efeitos como polarização térmica, perdas de calor, resistência ao transporte de massa e possíveis não idealidades experimentais. Esses desafios tornam-se ainda mais relevantes quando se passa de sistemas em escala de bancada para unidades em escala piloto, nas quais efeitos associados à geometria do módulo, distribuição de escoamento, instrumentação, perdas térmicas e variações operacionais podem reduzir a aderência de modelos simplificados às medições experimentais (\cite{Lisboa2024}).

% CITAÇÃO NECESSÁRIA: Verificar se ReviewAI_MD sustenta diretamente esta classificação em três categorias. Caso contrário, incluir uma referência adicional de revisão sobre modelagem em MD ou sobre modelagem físico-matemática, orientada por dados e híbrida.
De modo geral, as abordagens de modelagem aplicadas à destilação por membranas podem ser classificadas em três categorias principais: modelos baseados em primeiros princípios, modelos orientados por dados e abordagens híbridas (\cite{ReviewAI_MD}).

\subsection{Abordagens de modelagem física em destilação por membranas}

A modelagem física de sistemas de destilação por membranas pode ser desenvolvida com diferentes níveis de detalhamento, dependendo do objetivo do estudo, da configuração analisada e do custo computacional admissível. Em geral, essas abordagens buscam descrever os mecanismos acoplados de transferência de calor e massa que governam o processo, incluindo evaporação na interface quente, transporte de vapor através da membrana, difusão no domínio gasoso e condensação na região fria. Dessa forma, diferenciam-se de modelos puramente empíricos por partirem de leis físicas fundamentais, como balanços de massa e energia, equações de transporte e relações fenomenológicas associadas ao processo.

No contexto da destilação por membranas, uma primeira forma de modelagem física consiste na formulação de modelos distribuídos ou semi-distribuídos, capazes de representar variações espaciais das variáveis do processo. Esses modelos podem descrever gradientes de temperatura, concentração, pressão ou velocidade ao longo do módulo, permitindo uma análise mais detalhada dos fenômenos de transporte. No entanto, esse maior nível de detalhamento normalmente exige maior número de parâmetros, propriedades físicas, correlações de transporte e condições de contorno bem definidas.

Um exemplo dessa abordagem é o trabalho de \cite{Alsaadi2013}, que desenvolve um modelo unidimensional para módulos planos de AGMD. Nesse modelo, o módulo é discretizado ao longo da direção do escoamento e são consideradas as principais regiões do sistema: canal de alimentação aquecido, membrana polimérica, espaço de ar, filme de condensado, placa de resfriamento e canal frio. A formulação permite acompanhar a variação longitudinal das condições de operação e foi validada experimentalmente para avaliar o comportamento do processo em escoamentos co-corrente e contracorrente.

No modelo de \cite{Alsaadi2013}, são adotadas hipóteses como regime permanente, escoamento unidimensional, ausência de perdas de calor para o ambiente, pressão constante no espaço de ar e transporte de massa no espaço de ar predominantemente por difusão. O transporte de vapor através da membrana é representado pela combinação entre difusão molecular e difusão de Knudsen, enquanto a transferência de calor é descrita por mecanismos de convecção nos canais, condução através da membrana e do espaço de ar, e liberação de calor latente durante a condensação. Assim, o estudo exemplifica uma modelagem física mais detalhada, capaz de representar a distribuição longitudinal dos fenômenos de transporte, mas ainda baseada em hipóteses simplificadoras para manter o problema tratável.

Modelos físicos ainda mais detalhados podem ser formulados com base em abordagens bidimensionais e ferramentas de dinâmica dos fluidos computacional. Nesse sentido, \cite{Ansari2022} propõem um modelo bidimensional para AGMD em módulo plano contracorrente, com o objetivo de capturar a variação do fluxo de água ao longo do comprimento do módulo e comparar o desempenho da configuração AGMD com a DCMD. A formulação considera o canal de alimentação, a membrana hidrofóbica, o espaço de ar, o filme de condensado, a placa de resfriamento e o canal de permeado, sendo acoplada à solução das equações de continuidade, quantidade de movimento e energia por meio de CFD.

Essa abordagem permite representar de forma mais detalhada as variações espaciais do processo, sendo útil para investigar efeitos locais e limitações de desempenho que não seriam capturados por modelos concentrados. Por outro lado, a maior fidelidade física implica aumento da complexidade matemática e da demanda computacional, além de maior dependência de propriedades, correlações e condições de contorno. Assim, modelos distribuídos são particularmente úteis para análises mecanísticas detalhadas, mas podem ser menos práticos quando o objetivo é obter estimativas rápidas de desempenho ou integrar a modelagem física a rotinas de otimização e aprendizado de máquina.

Nesse contexto, modelos físicos reduzidos surgem como uma alternativa para representar o comportamento global do sistema com menor custo computacional. Diferentemente das formulações distribuídas ou bidimensionais, esses modelos utilizam grandezas médias ou efetivas, balanços globais e resistências equivalentes para estimar variáveis de desempenho. O modelo proposto por \cite{Lisboa2024}, por exemplo, representa uma formulação reduzida para um sistema V-AGMD em escala piloto, permitindo estimar grandezas como fluxo de permeado e temperaturas de saída sem resolver explicitamente os perfis espaciais ao longo do módulo.

A principal vantagem dos modelos reduzidos está na simplicidade de implementação e no menor custo de simulação, o que favorece sua utilização como referência física em pipelines de modelagem preditiva. Entretanto, essa simplificação também limita a representação de fenômenos locais, como gradientes axiais de temperatura, concentração e velocidade, efeitos hidrodinâmicos, perdas térmicas distribuídas e não idealidades geométricas do módulo. Dessa forma, modelos reduzidos apresentam um compromisso entre coerência física, baixo custo computacional e menor detalhamento fenomenológico.

A literatura mostra, portanto, que a modelagem física em destilação por membranas pode variar desde formulações unidimensionais ou bidimensionais, voltadas à descrição detalhada dos mecanismos de transporte, até modelos reduzidos de parâmetros concentrados, voltados à estimativa global do desempenho. Essa diversidade evidencia um compromisso entre fidelidade física, custo computacional, disponibilidade de parâmetros e aplicabilidade prática. No presente trabalho, esse panorama é relevante porque o modelo físico reduzido de \cite{Lisboa2024} é utilizado não como uma descrição completa do sistema, mas como uma referência físico-matemática simplificada, posteriormente incorporada à análise comparativa com modelos orientados por dados e abordagens híbridas.
\subsection{Modelos orientados por dados}

Modelos orientados por dados, também conhecidos como modelos \textit{data-driven}, surgem como uma alternativa atrativa quando se busca reduzir o custo computacional associado à simulação detalhada de sistemas de destilação por membranas. Em vez de resolver explicitamente as equações físicas do processo, esses modelos aprendem relações entre variáveis operacionais e saídas de interesse a partir de dados experimentais ou simulados, permitindo obter respostas rápidas após a etapa de treinamento.

% CITAÇÃO NECESSÁRIA: Inserir aqui trabalhos de MD que utilizaram planejamento experimental, regressão estatística ou metodologia de superfície de resposta, especialmente em unidades de bancada ou piloto.
Além das redes neurais artificiais, a literatura de MD também apresenta estudos baseados em planejamento experimental, regressão e metodologia de superfície de resposta. Nesses trabalhos, conjuntos planejados de experimentos são utilizados para ajustar equações empíricas capazes de relacionar variáveis operacionais, como temperatura, vazão, salinidade e pressão, a respostas do processo, como fluxo de permeado, rejeição de sais e indicadores de desempenho. Essas abordagens funcionam como metamodelos úteis para análise de sensibilidade e otimização local, especialmente em sistemas de bancada e unidades piloto.

% CITAÇÃO NECESSÁRIA: Verificar com a professora o trabalho/dissertação do Francisco, aluno da professora Carolina, que utilizou abordagem semelhante. Após localizar a referência, inserir aqui a contribuição específica do estudo.
No entanto, por serem ajustados a partir de um conjunto específico de dados, modelos empíricos baseados em planejamento experimental e regressão tendem a ser mais confiáveis dentro da faixa experimental considerada. Assim, embora sejam úteis para interpretação estatística dos efeitos operacionais e para otimização local, sua extrapolação para condições distintas daquelas utilizadas no ajuste deve ser realizada com cautela.

Na destilação por membranas, modelos orientados por dados também têm sido amplamente empregados na previsão do desempenho do processo por meio de redes neurais artificiais. Essas redes são utilizadas para mapear diretamente variáveis operacionais em saídas de interesse, como fluxo de permeado e temperaturas de saída (\cite{ANN_AGMD, ANN_VMD_Fouling}).

Esses modelos apresentam alta capacidade de aproximação de funções não lineares e bom desempenho dentro da faixa de treinamento. Contudo, sua generalização é limitada e dependente da distribuição dos dados (\cite{ReviewAI_MD}), podendo gerar previsões inconsistentes fora dessa região.

Dessa forma, embora os modelos orientados por dados apresentem potencial para representar relações não lineares em sistemas de MD, sua aplicação exige atenção à representatividade dos dados e à estratégia de validação adotada. Essa questão é particularmente relevante em sistemas experimentais organizados em regimes operacionais distintos, nos quais divisões aleatórias podem superestimar a capacidade de generalização do modelo.

\subsection{Abordagens híbridas}

Embora menos frequentes do que os modelos puramente físicos ou puramente orientados por dados, algumas abordagens buscam combinar conhecimento físico e aprendizado a partir de dados. Neste trabalho, essas estratégias são tratadas como abordagens híbridas, pois associam, em diferentes níveis, a coerência física dos modelos baseados em primeiros princípios à flexibilidade dos modelos ajustados a partir de dados experimentais ou simulados.

De modo geral, essas abordagens podem ser construídas de diferentes formas. Uma possibilidade consiste em incorporar equações físicas, balanços ou restrições do processo diretamente à formulação do modelo ou à função de perda. Outra possibilidade é utilizar um modelo físico como estimativa inicial e empregar um modelo orientado por dados para corrigir seus desvios em relação aos dados experimentais. Dessa forma, a informação física pode atuar tanto como elemento de regularização quanto como referência para o aprendizado de correções.

No contexto da destilação por membranas, uma das estratégias encontradas corresponde às \textit{Physics-Informed Neural Networks} (PINNs), nas quais as equações governantes são incorporadas à função de perda. O trabalho de \cite{PINN_MD}, por exemplo, utiliza esse tipo de abordagem para incluir conhecimento físico no treinamento da rede neural, buscando melhorar a consistência das previsões em relação aos fenômenos que regem o processo.

Outra classe de abordagem envolve modelos residuais, nos quais um modelo físico fornece uma estimativa inicial e uma rede neural é treinada para aprender a diferença entre essa estimativa e os dados experimentais. Essa estratégia é explorada em trabalhos como \cite{Zhou2024, zheng2021knowledge}, que destacam a possibilidade de combinar previsões baseadas em conhecimento físico com correções aprendidas a partir dos dados.

% CITAÇÃO NECESSÁRIA: Caso o trabalho do Francisco tenha utilizado abordagem híbrida, inserir aqui a referência e explicar a contribuição específica. Caso tenha utilizado planejamento experimental ou regressão, manter a referência na subseção de modelos orientados por dados.
As abordagens híbridas se inserem, portanto, como uma alternativa para combinar as vantagens dos modelos físicos e dos modelos orientados por dados. Sua principal contribuição está na tentativa de preservar coerência física sem abrir mão da capacidade preditiva dos métodos ajustados por dados, o que é particularmente relevante em sistemas experimentais complexos, com dados limitados ou sujeitos a não idealidades não completamente representadas por modelos simplificados.

\subsection{Síntese e lacunas}

A partir da literatura analisada, observa-se que a modelagem de sistemas de destilação por membranas tem sido conduzida por diferentes estratégias. Modelos baseados em primeiros princípios contribuem para a compreensão mecanística do processo e para a análise de condições operacionais não exploradas experimentalmente. Modelos empíricos, baseados em planejamento experimental e regressão, permitem construir metamodelos úteis para análise de sensibilidade e otimização local. Já os modelos orientados por dados, especialmente redes neurais artificiais, ampliam a capacidade de representar relações não lineares entre variáveis operacionais e respostas do processo. Por fim, as abordagens híbridas buscam combinar coerência física e capacidade preditiva, utilizando o conhecimento físico como restrição, referência ou ponto de partida para modelos ajustados a partir de dados.

Apesar desses avanços, a literatura ainda apresenta limitações metodológicas relevantes. Predomina o uso de divisões simples dos dados, como 70/15/15 ou 80/20, geralmente realizadas de forma aleatória (\cite{ANN_VMD_Fouling, ReviewAI_MD}). Essa prática desconsidera a estrutura dos dados experimentais, como regimes operacionais distintos e dependências entre amostras.

Além disso, raramente são utilizadas estratégias mais robustas de validação, como validação cruzada, sendo comum a avaliação baseada em uma única divisão dos dados (\cite{ANN_VMD_Fouling}), o que pode introduzir viés. Essa limitação é particularmente importante em sistemas experimentais nos quais diferentes amostras podem estar associadas a um mesmo regime operacional ou a condições próximas de operação.

A seleção de modelos e hiperparâmetros é frequentemente realizada de forma heurística, com otimização limitada e baseada apenas na minimização de métricas como RMSE ou MSE, sem consideração explícita de complexidade ou variabilidade estatística (\cite{ANN_VMD_Fouling}). Adicionalmente, a maioria dos estudos avalia os modelos apenas em regime de interpolação, com pouca análise da capacidade de extrapolação (\cite{ANN_VMD_Fouling, ReviewAI_MD}).

Alguns trabalhos, contudo, adotam abordagens mais robustas. O estudo de \cite{PINN_MD} utiliza validação cruzada e avalia cenários futuros, enquanto \cite{ANN_AGMD} emprega validação externa independente. Já \cite{Zhou2024} apresenta uma abordagem híbrida com busca sistemática de hiperparâmetros e análise explícita de generalização.

Dessa forma, evidencia-se uma lacuna na adoção de estratégias de validação mais rigorosas, na consideração da estrutura dos dados experimentais e na análise da capacidade de generalização dos modelos em sistemas reais. Nesse contexto, este trabalho contribui ao investigar modelos orientados por dados e abordagens híbridas aplicados a um sistema V-AGMD em escala piloto, incorporando informação física proveniente de um modelo reduzido e adotando uma estratégia de validação compatível com a organização dos dados experimentais.